\scanpage{28}

techniquement \unsure{simple|souple}. Si $X$ est \deleted{\unsure{paracompact}} \inserted{un} espace paracompact, on pourra prendre (à titre provisoire !) pour $\mathcal V_{Y,X}$ le ``germe \unsure{d'espaces|d'espace}'' de $X$ autour de $Y$ -- qui est donc un topos, sav. la $\varprojlim$ toposique des voisinages ouverts de $Y$ dans $X$. Intuitivement, $\mathcal V_{Y,X}\hookrightarrow X$ a les propriétés d'une imm. ouverte, et $Y\hookrightarrow\unsure{\mathcal V_{Y,X}|\deleted{\mathcal V}_{Y,X}}$ celles d'une immersion fermée, et (3) constitue une factorisation canonique de (1).

% The cancellation has a long top stroke, an open right side, a stepped lower-left edge and four diagonals; the words remain ordinary transcribed text.
\deleted{L'hypothèse de locale trivialité normale implique que l'on a (moyennant des conditions de ``modération'' convenables) une inclusion}

\margin{L'intérêt principal de cette factorisation, c'est que \unsure{l'inclusion}''
\begin{tikzcd}[ampersand replacement=\&,column sep=small,cramped]
Y\arrow[r,-,to path={(\tikztostart)--(\tikztotarget) node[circle,fill=black,inner sep=.45pt] {} \tikztonodes}]\&\mathcal V_{Y,X}
\end{tikzcd}
est une équivalence d'homotopie'' (indépendamment de toute hyp. de locale trivialité)}

(4)
\diagram{\begin{tikzcd}[ampersand replacement=\&]
Y\arrow[r,"\text{équiv. d'homotopie}"]\&\mathcal V_{Y,X}
\end{tikzcd}}

Supposons p. ex. $Y$ et $X-Y=X^*$ \deleted{des} variétés, et $Z\inserted{{}^*=Z\setminus\{z\}}$ une variété (moyennant des hyp. de modération, l'hyp. sur $Z$ devrait être une conséquence de la première hypothèse), \deleted{\unsure{$Z^*$ est|et que}} alors que $z$ est un pt singulier \inserted{de $Z$}, i.e. $X$ est ``singulier le long de $Y$''. Il est naturel alors de vouloir \unsure{réconstituer|reconstituer} $X$ en termes des variétés $Y$ et $X^*$,

\margin{27}
